\centerline{\bf \Huge Делимость}
\title{Делимость}
\begin{flushright}
    \scriptsize{2...3...5...11...простые числа - это числа,\\ которые делятся только на 1 и на себя... \\они дают мне силу.\\ Энрико Пуччи. Невероятные приключения ДжоДжо. Каменный океан.\\Том 6 глава 49}
\end{flushright}
Прежде чем говорить о делимости, нужно понять, а что значит, что одно число делится на другое. Что значит, что $a$ делится $b$? 

\textbf{Определение.} Число $a$ делится на число $b$, если $a$ можно представить, как $a=bn$, где $n$ это какое-то целое число. Еще говорят число $a$ кратно числу $b$. Обозначается как $a \ \vdots \ b$

Иными словами одно число делится на другое, если одно число какое-то количество раз входит в другое. Например, 6 делится на 2 и представляется, как $6 = 2 \cdot 3$. То есть двойка входит три раза в шестерку. Такое представление очень удобное при работе с числами, если мы знаем, на что они делятся. Например, числа делящиеся на 2 представляются в виде $2n$, числа делящиеся на 7 представляются в виде $7n$, делящиеся на 147 представляются в виде $147n$ и так далее.\\

\noindent\textbf{\Large Свойства делимости:}\\
1. Если $a$ делится на $b$, а $b$ в свою очередь делится на $c$, то $a$ делится на $c$.\\
2. Если $a$ делится на $c$ и $b$ делится на $c$, то $(a+b) \ \vdots \ c$ и $(a-b) \ \vdots \ c$\\
3. Если $a$ делится на $c$, но $b$ не делится на $c$, то $(a+b)$ и $(a-b)$ не делятся на $c$\\
4. \textbf{Очень важный факт!} Если $ab$ делится на $c$ и $(c,b)=1$ (эта запись означает, что их НОД равен единице), то $a$ \ \vdots \ $c$ \\

\noindent\textbf{\Large Признаки делимости:}\\ 
\\
\textbf{Признак делимости на 2.} Число делится на 2 тогда и только тогда, когда оно оканчивается на четную цифру;\\
\textbf{Признак делимости на 4.} Число делится на 4 тогда и только тогда, когда две последние его цифры образуют число, делящееся на 4;\\
\textbf{Признак делимости на 8.} Число делится на 8 тогда и только тогда, когда три последние его цифры образуют число, делящееся на 8; (что-то тут нечисто...)
\newpage
\noindent \textbf{Признак делимости на 3.} Число делится на 3 тогда и только тогда, когда сумма его цифр делится на 3;\\
\textbf{Признак делимости на 9.} Число делится на 9 тогда и только тогда, когда сумма цифр делится на 9.\\
\textbf{Признак делимости на 11.} Число делится на 11 тогда и только тогда, когда модуль разности между суммой цифр, занимающих нечётные позиции, и суммой цифр, занимающих чётные места, делится на 11.

Это основные признаки делимости, которые вам понадобятся для решения задач в дз, да и в целом на всех Олимпиадах. А теперь перейдем к задачам!\\

\noindent\textbf{\Large{Четность}}

\problem{1}Чётными или нечётными будут сумма и произведение:\\
а) двух чётных чисел?\\
б) двух нечётных чисел?\\
в) чётного и нечётного?\\

\textbf{Решение:} Для решения этой задачи нам как раз поможет представление чётного числа. Но как же представить нечётное число? Очень просто! Нечётное число это такое число, которое при делении на 2 даёт остаток 1, то есть представляется в виде $2k+1$. Решаем по порядку. В пункте а) просят узнать что будет с суммой и произведением двух чётных чисел. Возьмем два произвольных чётных числа $2m$ и $2n$. Сложим их и получим число $2m+2n = 2(m+n)$, которое очевидно чётное. Если взять их произведение, то получим $2m\cdot2n=4mn$, что тоже делится на 2. В пункте б) требуют сделать то же самое, но уже с нечетными числами. Возьмем сумму двух произвольных нечётных чисел $2m+1$ и $2n+1$. Получим $2m+1+2n+1=2m+2n+2=2(m+n+1)$, что делится на 2. Значит нечётное плюс нечётное будет чётное. Но их произведение $(2m+1)(2n+1)=4mn+2m+2n+1=2(mn+m+n)+1$ уже будет нечётным. Пункт в) сделайте сами :)

Некоторым продолжением этой задачи является следующая. 

\problem{2} Доказать, что сумма\\
а) любого количества чётных слагаемых чётна\\
б) чётного количества нечётных слагаемых чётна\\
в) нечётного количества нечётных слагаемых нечётна\\

\textbf{Решение:} С пунктом а) всё предельно понятно. Складывай чётные числа сколько хочешь, каждый раз будешь получать чётное. А вот с нечётными нужно следить за чётностью суммы. Если сложим два нечётных, то как мы уже знаем, получится чётное. Если прибавить к ним третье нечётное, получится опять нечётное. Если прибавить к ним чётвертое нечётное, то вновь получим чётное. Несложно прослеживается закономерность, что когда мы складываем чётное число нечётных чисел, то получим чётное, а нечётное число нечётных чисел, то нечётное.

\problem{3}Баир сложил два натуральных чисел, а потом умножил результат на произведение его чисел. У него получился ответ 2023. Докажите, что Баир опять все неправильно посчитал...

\textbf{Ответ: Нет, не могло.}

\textbf{Решение:} Анджур Аркадьевич уже научил нас вводить переменную на первом вебинаре. Давайте воспользуемся этим знанием! Пусть $a$ это первое число, $b$ это второе число. Пусть так получилось, что сумма этих чисел помноженная на произведение стала равна 2023. Тогда получим, что $ab(a+b)=2023$. Казалось бы страшное буквенное выражение, на которое даже смотреть не хочется. Но если знать, что такое чётность то и думать не придётся. Заметим, что ни $a$, ни $b$ не могут быть чётными, иначе левая часть будет чётной, а правая у нас нечётная. Тогда остается только вариант, когда они оба нечётные. Но слева мы умножаем на сумму этих чисел. Мы знаем, что сумма двух нечётных чётна, а значит умножаем мы слева на чётное число, то есть слева число опять чётное, а справа нет. Противоречие! Значит такого быть не могло.\\

\noindent\textbf{\Large Признаки делимости}

\problem{1} Делится ли число $10^{2024}+8$ на $9$?

\textbf{Ответ: Да, делится.}

\textbf{Решение:} Число $10^{2024}$ это просто единичка и много ноликов. Если прибавить к этому числу 8, то получится единичка, все так же куча ноликов и в конце 8. Сумма цифр такого числа равна 9, а значит число делится на 9 по признаку делимости.

\problem{2}\textit{Бессмертная классика.} РиЧо и ДжурАн хотят договориться о встрече и РиЧо послал ДжуРану в телеграме свой номер телефона. Но канал связи был плохой и одна цифра не успела передаться.
\newpage
\noindent В итоге ДжурАну пришло вот такое число 10333147966386144929\textunderscore66651337523200000000. РиЧо успел передать, что его число совпадает с числом 35! (35 факториал). Помогите ДжурАну восстановить номер телефона РиЧо, чтобы они смогли договориться о встрече и сыграть в БолФут.

\textbf{Ответ: 6.}

\textbf{Решение:} Поскольку 35! делится на 9, сумма цифр этого числа также делится на 9. Нетрудно подсчитать, что сумма цифр (за исключением пробела) написанного числа даёт остаток 3 при делении на 9. Отсюда следует, что цифра, стоявшая на месте пробела, равна 6.


\problem{3} Злой Расул Владимирович в качестве наказания за плохую арифметику решил забрать у Баира все его скины в Standoff 2, но тот поставил на телефон семизначный пароль. Сумеет ли Расул Владимирович наказать Баира, если он знает, что пароль состоит только из двоек и троек, делится на 12 и двоек там больше чем троек? 

\textbf{Ответ: Да, сможет.}

\textbf{Решение:} В силу признака делимости на 4 код может оканчиваться только цифрами 32 (другие двузначные числа, составленные из цифр 2 и 3, не делятся на 4).

Двоек в коде больше, чем троек; значит, двоек не меньше четырёх, а троек не больше трёх. Если в коде четыре двойки и три тройки, то сумма цифр кода равна 2 · 4 + 3 · 3 = 17 и не делится на 3, поэтому и сам код не делится на 3. По аналогичной причине код не может состоять из пяти двоек и двух троек (тогда сумма цифр была бы равна 2 · 5 + 3 · 2 = 16). Значит, код может состоять только из одной тройки и шести двоек (тогда сумма цифр равна 2 · 6 + 3 · 1 = 15 и код делится на 3).

Положение единственной тройки в коде мы уже определили, а остальные цифры · двойки. Значит, подходит только код 2222232.\\ 

\noindent\textbf{\Large Разложение числа по степеням десятки}

Все мы знаем, что в каждом числе есть разряд единиц, разряд десятков, сотен, тысяч и так далее. Но как эти названия связаны с самим числом? Почему наши числа зачастую называют \textit{десятичными}? Все на самом деле очень просто. Это такой конструктор, из которого мы собираем числа. Деталями этого конструктора служат степени десятки. Давайте посмотрим, как это работает на деле. Посмотрим на число 
\newpage
\noindent 4325. Как нам получить последнюю цифру? Нужно 5 раз взять единицу. Как к этой единице слева еще приписать двойку? Нужно 2 раза взять десятку. Потом 3 раза сотню и 4 раза тысячу. Таким образом $4325 = 4\cdot1000+3\cdot100+2\cdot10+5\cdot1$. И так можно представить любое число! Это и есть волшебное разложение числа по степеням десятки. В общем виде это выглядит вот так:
Пусть $\overline{abc}$ произвольное трехзначное число, где a,b и с это его цифры. Тогда оно представляется в виде $\overline{abc}=100a+10b+c$. Зачем же это нужно? Сейчас увидим в задачах.

\problem{1}Найдите все двузначные числа, равные удвоенной сумме своих цифр.

\textbf{Ответ: 18.}

\textbf{Решение:} Двузначное число имеет 
 вид $\overline{ab}=10a+b$, где $a$ и $b$ его цифры. Тогда сумма его цифр это $a+b$. По условию нам говорят, что $10a+b=2(a+b)$. Раскроем скобки и сократим подобные слагаемые и получим, что $8a=b$. Вспоминим, что $a$ и $b$ это \textbf{цифры}, а они у нас бывают от 0 до 9. Значит $a$ не может быть сильно большим, даже двойке быть равным уже не может. Значит $a=1, b=8$. То есть наше число это 18. А перебором такая задача делалась бы очень долго!

 \problem{2}Тимур задумал три различные цифры, отличные от нуля. Эрдем записал все возможные двузначные числа, в десятичной записи которых использовались только эти цифры. Сумма записанных чисел равна 198. Найдите цифры, задуманные Тимуром
 
 \textbf{Ответ: 1, 2, 3.}

 \textbf{Решение:} Пусть а, b и с — три цифры, задуманные Тимуром. Существует девять двузначных чисел, в десятичной записи которых используются только эти цифры: $ \overline{aa}$; $ \overline{bb}$; $ \overline{cc}$; $ \overline{ab}$; $ \overline{ba}$; $ \overline{ac}$; $ \overline{ca}$; $ \overline{bc}$; $ \overline{cb}$. Найдем их сумму, разложив каждое из чисел в виде суммы разрядных слагаемых: (10a + a) + (10b + b) + (10c + c) + (10a + b) + (10b + a) + (10a + c) + (10c + a) + (10b + c) + (10c + b) = 33a + 33b + 33c = 33(a + b + c). По условию, 33(a + b + c) = 198, то есть, a + b + c = 6. Существует единственная тройка различных и отличных от нуля цифр, сумма которых равна 6. Это 1, 2 и 3.
 
 На самом деле эту задачу можно было решить ещё легче, зная комбинаторику. Как раз на вебинаре по комбе мы вернемся к этой задаче и приведем альтернативное решение.

 \textbf{\Large Теория чисел в играх}

 \problem{1}Чтобы у героя Axe из Dota 2 сработала пассивная способность counter helix нужно 4 раза по нему ударить. Во время игры на него напало $3^{2023}$ противников и каждый ударил ровно по одному разу. Сколько раз по нему нужно ещё ударить, чтобы ещё раз сработал counter helix?

 \textbf{Решение:} Это классическая задача на нахождение остатка при делении на какое-то число. Axe крутится каждый четвертый удар по нему. Значит нам нужно разделить число $3^{2023}$ на 4 и узнать остаток. Вручную это делать, конечно, не стотит, поэтому мы будем хитрить. Начнем с малого. Какой остаток дает 3 при делении на 4? Конечно же 3. Какой остаток дает $3^2$ при делении на 4. Девятка дает остаток 1. Идем дальше $3^3$ дает остаток 3, $3^4$ дает 3, потом опять 1, опять 3 и так далее. В чётных степенях тройка дает остаток 1, а в нечётных 3. Стало быть наше число даёт остаток 3 при делении на 4, а значит осталось ударить ещё один раз, чтобы сработала способность. По-хорошему нужно аккуратно показать, почему остатки себя так ведут, но в рамках наших занятий мы этого не успеем:(

\problem{2} Очир и Расул проводят бой в Hearthstone. У Очира на поле боя есть карта Принц Кель'тас, которая позволяет ему применять каждое третье заклинание бесплатно. Чтобы увеличить шансы на победу, Очир включил читы и применил сразу $2023^{2023}$ заклинаний. Сколько ему ещё нужно применить заклинаний, чтобы опять сработал эффект Принца Кель'таса 

 \begin{figure}[h] 
    \hspace*{-12cm}\includegraphics[scale=0.12]{ЭЛМШ_2/Конспекты занятий/prince.png}
\end{figure}

\textbf{Ответ: 2 раза.}

\textbf{Решение:} Эта задача очень похожа на предыдущую, но нюанс в том, что как-то не хочется возводить 2023 даже во вторую степень. На помощь приходят сравнения по
\newpage
\noindent модулю. Очень мощный инструмент в Олимпиадной теории чисел. До 8 класса многие задачи убиваются именно с помощью этого инструмента. Овладев им, вам откроются большие просто в теории чисел. Дело в том, что если мы хотим узнать остаток $2023^{2024}$ при делении на 3, то нам не обязательно смотреть брать 2023, а только лишь его остаток при делении на 3. То есть вместо $2023^{2024}$ смотреть на $1^{2024}$. А тут задача становится уж совсем очевидной. Получаем, что $2023^{2024}$ дает остаток 1 при делении на 3, а значит осталось применить еще 2 заклинания.